\section{Из чего складывается хроматичность}
\label{sec:terms}

Разложение гамильтониана \eqref{eq:H} около дисперсионной орбиты даёт $\mathcal{H}_2^{(1)}$ в виде суммы нескольких слагаемых разного происхождения. Каждое из них после усреднения \eqref{eq:avg} даёт свой вклад в $I_{x,y}$ и в $\xi_{x,y}$. В этой главе перечислены вклады тела элементов: кинематический (он же квадрупольный), кривизна тела диполя, перекрёстный член с $D_1'$ и секступоли. Вклад края диполя отложен до главы~\ref{sec:edge}.

\subsection{Разложение гамильтониана}

До третьего порядка по $(x,p_x,p_y,\delta)$ корень в \eqref{eq:H} равен $(1+\delta)-(p_x^2+p_y^2)/(2(1+\delta))$, и
\begin{equation}
\label{eq:H3}
\mathcal{H}\simeq
-\frac{x\,\delta}{\rho}
+\Bigl(1+\frac{x}{\rho}\Bigr)\frac{p_x^2+p_y^2}{2(1+\delta)}
+\frac{K_x}{2}x^2+\frac{K_y}{2}y^2
+\frac{K_2}{6}\bigl(x^3-3xy^2\bigr)
\end{equation}
(без постоянной). Первый член --- источник дисперсии, $K_xx^2/2$ и $K_yy^2/2$ от $\delta$ не зависят. Подстановка $x=D_1\delta+X$, $p_x=D_1'\delta+P_x$ и сбор членов $\delta\times(\text{квадратичные по }X,P_x,Y,P_y)$ дают
\begin{equation}
\label{eq:H21}
\mathcal{H}_2^{(1)}
=\underbrace{-\frac{P_x^2+P_y^2}{2}}_{\mathrm{A}}
+\underbrace{\frac{D_1}{\rho}\,\frac{P_x^2+P_y^2}{2}}_{\mathrm{C}}
+\underbrace{\frac{D_1'}{\rho}\,X P_x}_{\mathrm{E}}
+\underbrace{\frac{K_2D_1}{2}\bigl(X^2-Y^2\bigr)}_{\mathrm{S}}.
\end{equation}
Член A --- от множителя $1/(1+\delta)$, C --- от множителя $1+x/\rho$ при сдвиге $x$ на $D_1\delta$, E --- от того же множителя при сдвиге $p_x$ на $D_1'\delta$ (произведение $X/\rho$ на $2D_1'\delta P_x$), S --- от кубического члена секступоля.

\subsection{Вклады по членам}

По \eqref{eq:xi-ham}, \eqref{eq:avg}
\begin{equation}
\label{eq:I-all}
I_x=-\gamma_x+\frac{D_1}{\rho}\gamma_x-\frac{2D_1'}{\rho}\alpha_x+K_2D_1\beta_x,
\qquad
I_y=-\gamma_y+\frac{D_1}{\rho}\gamma_y-K_2D_1\beta_y,
\end{equation}
что совпадает с выражениями~\cite{EqNote}: $I_x=-\gamma_x(1-D_1/\rho)-(2\alpha_x/\rho)D_1'$, $I_y=-\gamma_y(1-D_1/\rho)$ в диполе и $I=-\gamma\pm K_2D_1\beta$ на квадруполях, секступолях и промежутках. Отдельно:
\begin{align}
\label{eq:xiA}
\xi^{\mathrm{A}}_{x,y}&=-\frac{1}{4\pi}\oint\gamma_{x,y}\,\mathrm{d}s
=-\frac{1}{4\pi}\oint K_{x,y}\beta_{x,y}\,\mathrm{d}s,\\
\label{eq:xiC}
\xi^{\mathrm{C}}_{x,y}&=+\frac{1}{4\pi}\int_{\text{диполи}}\frac{D_1\gamma_{x,y}}{\rho}\,\mathrm{d}s,\\
\label{eq:xiE}
\xi^{\mathrm{E}}_{x}&=-\frac{1}{2\pi}\int_{\text{диполи}}\frac{D_1'\alpha_x}{\rho}\,\mathrm{d}s,
\qquad\xi^{\mathrm{E}}_y=0,\\
\label{eq:xiS}
\xi^{\mathrm{S}}_{x,y}&=\pm\frac{1}{4\pi}\oint K_2D_1\beta_{x,y}\,\mathrm{d}s.
\end{align}
Во второй форме \eqref{eq:xiA} использовано $K\beta=\gamma+\alpha'$ (следует из уравнения для бетатронной функции) и $\oint\alpha'\,\mathrm{d}s=0$. Эта замена законна только для интеграла по всему кольцу, причём $K_x$ включает слабую фокусировку $1/\rho^2$ диполей.

\paragraph{Член A: кинематический.}
Он есть в любом элементе, в том числе в дрейфе: в канонических переменных фокусировка от $\delta$ не зависит, а скорость поворота $x'=p_x/(1+\delta)$ --- зависит. В форме $-\oint K\beta/4\pi$ он выглядит как <<квадрупольная>> хроматичность $\Delta K=-K\delta$ уравнения \eqref{eq:hill}, но происходит из кинетического члена. По этой причине формулу $\Delta\nu=(1/4\pi)\oint\beta\Delta K$ нельзя применять поэлементно внутри диполя~\cite{EqNote}: там $x'\neq p_x$ уже в линейном по $\delta$ порядке, и члены C, E не сводятся к $\Delta K$.

\paragraph{Член C: кривизна тела.}
Частица на дисперсионной орбите в диполе идёт по большему ($D_1>0$) радиусу, и её путь на единицу $s$ длиннее в $1+D_1\delta/\rho$ раз. Бетатронный набег фазы растёт пропорционально; член одинаков по форме для $x$ и $y$ (в $y$ --- через $\gamma_y$).

\paragraph{Член E: перекрёстный.}
Наклон дисперсионной орбиты $D_1'\delta$ вместе с множителем $1+X/\rho$ даёт член $XP_x$, т.\,е. связь координаты и импульса горизонтальных колебаний. Он только горизонтальный: аналогичный член по $y$ требовал бы вертикальной дисперсии. Его среднее $\propto\alpha_x$, поэтому член чувствителен к тому, где в диполе находится перетяжка $\beta_x$.

\paragraph{Член S: секступоли.}
Секступоль на дисперсионной орбите --- квадруполь с градиентом $K_2D_1\delta$. Знак разный для $x$ и $y$; это основа коррекции хроматичности двумя семействами~\cite{ValetovPhD}.

\paragraph{Сумма тела.}
Сумма \eqref{eq:xiA}--\eqref{eq:xiS} есть хроматичность гамильтониана \eqref{eq:H}, в котором поле диполя обрывается на торцах без продольной компоненты. Такая модель <<только тело>> точна внутри элементов, но поле с разрывом $B_y$ и $B_s=0$ не удовлетворяет $\operatorname{rot}\mathbf{B}=0$ на границе. Недостающий вклад --- предмет следующей главы.
